\documentclass[11pt]{article}
\usepackage{amsmath,amssymb,amsthm}
\usepackage[margin=1in]{geometry}
\usepackage{hyperref}

\newtheorem{theorem}{Theorem}
\newtheorem{lemma}[theorem]{Lemma}
\newtheorem{proposition}[theorem]{Proposition}
\newtheorem{remark}[theorem]{Remark}
\newtheorem{definition}[theorem]{Definition}
\newtheorem{conjecture}[theorem]{Conjecture}

\newcommand{\zD}{\zeta_D}
\newcommand{\PE}{\mathrm{PE}}

\title{Term-by-Term Extraction of Euler Products into\\ Ratios of Riemann Zeta Values\\
\large Working notes, with Mathar (arXiv:1106.4038) as prior art}
\author{}
\date{\today}

\begin{document}
\maketitle

\begin{abstract}
For a multiplicative arithmetic function $a(n)$ with local (Bell series) factor
$B_p(x) = \sum_{e\ge0} a(p^e)x^e$ at $x=p^{-s}$, we consider the problem of writing
the Euler product $\zD(s)=\prod_p B_p(p^{-s})$ as an (in general infinite) product
of Riemann zeta values $\zeta(u_is-l_i)^{\pm\gamma_i}$. We record the exact
statement of the underlying identity, formalize the term-by-term extraction step
as an algorithm on exact rational functions of $x=p^{-s}$, and identify
R.~J.~Mathar's survey \cite{Mathar2011} as an existing treatment of the same
question, with a worked example (the Hadamard square of Euler's totient) that is
structurally identical to our target case $1/\varphi(n)^2$. We use this to
calibrate our own algorithm before pushing it further.
\end{abstract}

\tableofcontents

\section{Setup}

\begin{definition}[Dirichlet series]
For an arithmetic function $a:\mathbb{Z}_{\ge1}\to\mathbb{C}$,
\[
\zD(s) \;\equiv\; \sum_{n\ge1} \frac{a(n)}{n^s}.
\]
\end{definition}

If $a$ is multiplicative, $\zD$ factors over primes into local Bell series
\[
\zD(s) = \prod_p B_p(p^{-s}), \qquad B_p(x) \equiv \sum_{e\ge0} a(p^e)\, x^e .
\]
For all functions of interest here $B_p(x)$ is a \emph{rational} function of
$x$ with coefficients that are (typically polynomial or rational) functions of
$p$ — this rationality, not multiplicativity per se, is what makes the zeta
extraction go through, since it lets us write $B_p(x) = \mathrm{NUM}_p(x)/\mathrm{DEN}_p(x)$
with $\mathrm{NUM}_p,\mathrm{DEN}_p \in \mathbb{Q}(p)[x]$.

\subsection{The target identity}

Euler's product formula gives, for each fixed pair $(l,u)\in\mathbb{Z}\times\mathbb{Z}_{\ge1}$
and $\gamma\in\mathbb{Z}$,
\begin{align}
\prod_p (1 - p^{\,l-us})^{\gamma} &= \zeta(us-l)^{-\gamma}, \label{eq:euler-minus}\\
\prod_p (1 + p^{\,l-us})^{\gamma} &= \frac{\zeta(us-l)^{\gamma}}{\zeta(2us-2l)^{\gamma}}. \label{eq:euler-plus}
\end{align}
Equation \eqref{eq:euler-minus} is immediate from the Euler product of $\zeta$.
Equation \eqref{eq:euler-plus} follows by writing $1+y = (1-y^2)/(1-y)$ with
$y=p^{l-us}$ and applying \eqref{eq:euler-minus} twice.

The extraction problem is: given $B_p(x) = \mathrm{NUM}_p(x)/\mathrm{DEN}_p(x)$,
find a (possibly infinite) factorization
\begin{equation}
\label{eq:master-factorization}
B_p(x) = \prod_{i\ge1} (1 - S_i\, p^{\,l_i - u_i s})^{\gamma_i}, \qquad S_i\in\{+1,-1\},\ l_i,u_i,\gamma_i\in\mathbb{Z},\ u_i\ge1,
\end{equation}
ordered by increasing $u_i$ (and, for ties, decreasing $l_i$) so that the
factors are listed in order of decreasingly restrictive abscissa of convergence.
Substituting \eqref{eq:euler-minus}--\eqref{eq:euler-plus} termwise then gives
$\zD(s)$ as a product/ratio of zeta values.

\subsection{The extraction algorithm (informal statement)}

Our Mathematica prototype (\texttt{Extractor}/\texttt{GetCharsTerm2}/\texttt{Handler})
implements \eqref{eq:master-factorization} by:
\begin{enumerate}
\item Expand $\mathrm{NUM}_p(x)$ and $\mathrm{DEN}_p(x)$ as power series in $x$ to some working order $N$.
\item Identify the lowest-order term at which $\mathrm{NUM}_p$ and $\mathrm{DEN}_p$
(as formal series, normalized to constant term $1$) disagree; this fixes a
candidate binomial factor $(1-Sp^{l-us})^{\gamma}$ whose own series expansion
cancels the discrepancy to that order.
\item Divide it out (equivalently, multiply/divide $\mathrm{NUM}_p/\mathrm{DEN}_p$ by
$(1-Sp^{l-us})^{-\gamma}$) and recurse on the residual series.
\item Terminate if the residual becomes exactly $1$ (finite closed form); otherwise
continue and only conjecture the pattern in $(l_i,u_i,\gamma_i)$ from a run of
consecutive terms.
\end{enumerate}
This is a discretized, order-by-order relative of two established procedures:

\begin{itemize}
\item \textbf{Selberg--Delange.} Writing $F(s)=\zeta(s)^{\rho}G(s)$ with $G$
holomorphic and non-vanishing near $s=1$, only finitely many derivatives of
$G$ at $s=1$ are needed for the asymptotic count of $\sum_{n\le x}a(n)$. Our
procedure is more ambitious: we attempt to represent $G(s)$ \emph{itself} as an
infinite product of zeta values, which is a much stronger (and not always
achievable in closed form) statement.
\item \textbf{Plethystic exponential/logarithm} (Feng--Hanany--He and related
physics literature). Peeling off the next-order term of a generating function
and recursing is exactly the plethystic logarithm's algorithm; termination
occurs in the same ``nice'' (free / totally multiplicative generator) cases and
fails to terminate in general.
\end{itemize}

\section{Mathar's survey: the same construction, independently arrived at}

R.~J.~Mathar's survey \emph{Survey of Dirichlet Series of Multiplicative
Arithmetic Functions} (arXiv:1106.4038, published version cross-checked here)
\cite{Mathar2011} sets up \emph{precisely} the identity of
\S1.1 above, in his eq.~(1.5)--(1.7). Quoting the structure (paraphrased, our
numbering matches his where noted):

\begin{itemize}
\item His eq.~(1.5)--(1.6): the Bell series is written as an infinite product
\[
\zD(s) = \prod_p \prod_{i=1}^{\infty} (1-S_i p^{\,l_i-u_is})^{\gamma_i},
\]
with the ordering convention: smallest $u_i$ first, and for ties, largest
$l_i$ first — this is verbatim our ordering convention above, so there is no
ambiguity that this is the same object.
\item His eq.~(1.7) is exactly our \eqref{eq:euler-minus}--\eqref{eq:euler-plus}.
\item He is explicit that truncating this product is a numerical evaluation
tool for $\zD(s)$, and that in the ``nice'' cases the product over $p$ is a
\emph{finite} product/ratio of cyclotomic polynomials in $p^{-s}$, which then
becomes a finite product of zeta values; in the general case the expansion is
genuinely infinite.
\end{itemize}

So: Mathar's survey is not merely adjacent work, it is (as far as the general
framework goes) the same construction, published in 2011/2012. This does not
make our project redundant — his paper is a table of results with sketched
derivations, not a general algorithm with a reusable implementation, and he does
not treat the plethystic/Selberg--Delange framing at all — but it is the correct
reference point to (a) validate our extraction algorithm against, and (b) cite
as prior art.

\subsection{Worked example that matches our target: $\varphi^2(n)$}

Mathar's \S3.13.1 (his eq.~(3.75)--(3.79)) works the case most structurally
relevant to us. Euler's totient has master equation
\[
\varphi(p^e) = (p-1)p^{e-1}, \quad e>0,
\]
Bell series $\varphi(n)\mapsto \zeta(s-1)/\zeta(s)$ (his eq.~(3.76), a finite
closed form — the ``nice''/terminating case). Squaring the master equation
pointwise (Hadamard square, \emph{not} Dirichlet self-convolution) gives
\[
a(p^e) = (p-1)^2 p^{2e-2}, \qquad B_p(x) = 1+\sum_{e\ge1}(p-1)^2p^{2e-2}x^e
= \frac{1-2p^{1-s}+p^{-s}}{1-p^{2-s}}
= \frac{1-(2p-1)x}{1-p^2x}, \quad x=p^{-s}.
\]
This is his eq.~(3.78). The numerator is \emph{not} a perfect binomial power in
$x$ alone (the coefficient $2p-1$ is not a pure power of $p$), so the extraction
does not terminate at one factor; Mathar reports the result as an infinite
product (his eq.~(3.79)):
\begin{align}
\varphi^2(n) &\mapsto \zeta(s-2)\prod_p
(1-p^{1-s})^2(1+p^{-s})^{-1}(1-p^{2-2s})^{-1}(1+p^{1-2s})^{-2}\notag\\
&\quad\times(1-p^{3-3s})^{-2}(1+p^{2-3s})^{-4}(1-p^{1-3s})^{-2}
(1-p^{4-4s})^{-3}(1+p^{3-4s})^{-8}(1-p^{2-4s})^{-5}\notag\\
&\quad\times(1+p^{1-4s})^{-2}(1-p^{5-5s})^{-6}(1+p^{4-5s})^{-16}(1-p^{3-5s})^{-16}
(1+p^{2-5s})^{-8}(1-p^{1-5s})^{-2}\cdots,\ s>3. \label{eq:mathar-phi2}
\end{align}
The exponents $2,1,1,2,2,4,2,3,8,5,2,6,16,16,8,2,\dots$ are exactly the sort of
sequence our \texttt{Handler} step is meant to conjecture and (ideally) prove
closed-form recursions for; this gives us a fully independent, already-published
check case with 16 confirmed terms to validate our Haskell reimplementation
against before trusting it on new functions.

\begin{remark}
Note $\varphi^2(n)$ (pointwise square) is \emph{not} our target $1/\varphi(n)^2$
(pointwise reciprocal square). The two are related only in that both fail to be
expressible via the algebra of Dirichlet convolution applied to $\zeta$-ratios
directly (unlike, e.g., $\varphi(n)\star 1 = n$, which is transparent from
\eqref{eq:master-factorization} with a single finite factor). We work out our
own case in \S3.
\end{remark}

\section{Setting up $1/\varphi(n)^2$}

\begin{proposition}
\label{prop:bell-1overphi2}
The Bell series of $1/\varphi(n)^2$ at prime $p$ is
\[
B_p(x) = \frac{(p-1)^2p^2 + (2p-1)x}{(p-1)^2(p^2-x)}, \qquad x = p^{-s}.
\]
\end{proposition}

\begin{proof}
Master equation: $a(p^e) = 1/\varphi(p^e)^2 = (p-1)^{-2}p^{-2e+2}$ for $e\ge1$,
$a(1)=1$. Then
\[
B_p(x) = 1 + \sum_{e\ge1} \frac{p^{2-2e}}{(p-1)^2}\,x^e
= 1 + \frac{p^2}{(p-1)^2}\sum_{e\ge1}\Big(\frac{x}{p^2}\Big)^e
= 1 + \frac{p^2}{(p-1)^2}\cdot\frac{x/p^2}{1-x/p^2}
= 1 + \frac{x}{(p-1)^2}\cdot\frac{1}{1-x/p^2},
\]
valid for $|x|<p^2$ (i.e. $\Re(s)$ large enough, uniformly boundable since we
only need this as a formal identity of rational functions). Clearing
denominators,
\[
B_p(x) = \frac{(p-1)^2(p^2-x) + xp^2}{(p-1)^2(p^2-x)}
= \frac{(p-1)^2p^2 + x\big[p^2-(p-1)^2\big]}{(p-1)^2(p^2-x)}
= \frac{(p-1)^2p^2 + (2p-1)x}{(p-1)^2(p^2-x)},
\]
using $p^2-(p-1)^2 = 2p-1$.
\end{proof}

\subsection{Normalized form and first observations}

Dividing numerator and denominator by $(p-1)^2p^2$ to put $B_p(0)=1$:
\begin{equation}
\label{eq:phi-inv-sq-normalized}
B_p(x) = \frac{1+\dfrac{2p-1}{(p-1)^2p^2}\,x}{1-\dfrac{1}{p^2}x}.
\end{equation}

\begin{remark}[Flag: does not obviously match the algorithm's assumed template]
The denominator factor $1-p^{-2}x = 1-p^{-2-s}$ is exactly of the form
$1-p^{l-us}$ with $(l,u)=(-2,1)$, so by \eqref{eq:euler-minus} it contributes
$\zeta(s+2)^{-1}$ to $\zD$ directly — a clean, finite piece.

The numerator, however, has coefficient $\dfrac{2p-1}{(p-1)^2p^2}$, which is
\emph{not} a Laurent monomial in $p$ (it is a proper rational function of $p$
with a pole at $p=1$, which never occurs among primes but is still not of the
$p^{c}$ shape our extraction step pattern-matches on term by term). This means:
\begin{itemize}
\item Either the extraction must proceed by first expanding this coefficient
in a partial-fraction / asymptotic series in $1/p$ (i.e. $\frac{2p-1}{(p-1)^2p^2}
= \sum_k c_k p^{-k}$ for appropriate rational $c_k$, valid for the range of $p$
that matters, and each such term then seeds its own zeta factor at the
appropriate order) — effectively the numerator itself is not degree-1 in $x$
once expanded this way, it is a full power series in $x$ whose $p$-dependence
at each order in $x$ must independently be decomposed into $p$-power monomials;
\item or there is no finite factorization and the extraction is a genuinely
infinite, non-terminating recursion whose $\gamma_i$ pattern we would need to
conjecture numerically (as Mathar did for $\varphi^2(n)$) rather than derive
in closed form.
\end{itemize}
This needs to be checked computationally (run the term-by-term extractor to
enough orders and look for a recognizable exponent sequence, ideally cross-checked
against OEIS the way Mathar does) before we assert either alternative. Flagging
this explicitly rather than asserting a closed form we have not verified.
\end{remark}

\section{Next steps}

\begin{enumerate}
\item Reimplement the extraction algorithm in Haskell on exact \texttt{Rational}
coefficients, first validated against Mathar's $\varphi^2(n)$ example
\eqref{eq:mathar-phi2} (16 known terms) and against the trivially-finite Jordan
totient case $J_k(n) \mapsto \zeta(s-k)/\zeta(s)$ (his eq.~(3.87)) as a
terminating sanity check.
\item Once validated, run it on Proposition~\ref{prop:bell-1overphi2}'s $B_p(x)$
for $1/\varphi(n)^2$ to sufficient order in $x$ and extract/conjecture the
$(l_i,u_i,\gamma_i)$ sequence.
\item Search OEIS for the resulting exponent sequence, the way Mathar
cross-references A-numbers throughout — this is a cheap and effective
correctness check.
\item Only after (2)--(3) succeed, write the closed-form (or asymptotic infinite
product) statement for $1/\varphi(n)^2$ as a theorem with proof, in the same
style as Proposition~\ref{prop:bell-1overphi2}.
\end{enumerate}

\begin{thebibliography}{9}
\bibitem{Mathar2011}
R.~J.~Mathar, \emph{Survey of Dirichlet Series of Multiplicative Arithmetic
Functions}, arXiv:1106.4038 [math.NT].
\end{thebibliography}

\end{document}